\section{The actual Masullo kernel: a diagonal bound, an indentation defect, and a concrete repair}
\label{sec:SM-actual-kernel}

We use the operators actually named in the manuscript:
\[
 H=W_I^s(C_F)=L^2(C_F,w_{I,s}\,d\mu),\quad
 S_\phi=A(\phi)P_a^{\mathrm{spec}},\quad
 T_{\rho,\phi}=M_\rho A(\phi)P_a^{\mathrm{geom}}.
\]
The projections are Definition~10.4.1.1, p.~456; nuclearity for
$S_\phi$ is asserted in Theorem~10.4.2.8, pp.~463--464.
Definition~10.4.3.7, pp.~473--474, gives
$k^{\mathrm{geom}}_\phi(x,y)=\phi(xy^{-1})-k^{\mathrm{spec}}_\phi(x,y)$.
Equation~(10.4.3.48), p.~473, assigns to $T_{\rho,\phi}$ the
weighted-measure kernel
$w_{I,s}(y)^{-1}\rho(x)k^{\mathrm{geom}}_\phi(x,y)$.
The first result is conditional on the indicated actual $S_\phi$
being trace class and on this kernel identification.  It does not
replace that operator by a separately constructed zero spectrum or
recheck the entire earlier nuclearity argument.

\subsection{An explicit diagonal convergence estimate}

\paragraph{Finding F-SM-18 (trace-density domination and valid-indentation shell convergence).}
Let $\nu=w\mu$, where $\mu$ is a positive Radon measure and $w>0$.
Every trace-class operator $S$ on $L^2(X,\nu)$ has a unique
$d_S^\nu\in L^1(X,\nu)$ with
\[
 \operatorname{Tr}(M_fS)=\int_X f\,d_S^\nu\,d\nu
 \quad(f\in L^\infty),\qquad \|d_S^\nu\|_1\leq\|S\|_1.
 \tag{SM.18.1}
\]
Its density relative to $\mu$ is $d_S^\mu=w\,d_S^\nu$, with
the same $L^1$ norm.  For the actual $S_\phi$, if $\phi(1)=0$,
the trace-density geometric diagonal is $d=-d_{S_\phi}^\mu$, so
\[
 \|d\|_{L^1(\mu)}\leq\|S_\phi\|_1.
 \tag{SM.18.2}
\]
When the manuscript's averaging procedure gives the trace
diagonal of its spectral kernel, this agrees with its smoothened
geometric diagonal, since the convolution kernel has continuous
diagonal $\phi(1)$.  Its cutoff diagonal is $\rho d$;
for a sequence $|\rho_j|\leq1$ with $\rho_j\to1$ almost
everywhere, $\rho_jd\to d$ in $L^1$.

Suppose now $X=G/H$, with $G$ second countable, locally compact
abelian and $H$ closed and discrete with counting measure.
Take $\mu$ to be the quotient Haar measure and $D$ a measurable
fundamental domain carrying its corresponding measure $\mu_D$.
For a nonnegative measurable normalized indentation $b$, put
\[
 \sum_{r\in H}b(r^{-1}x)=1,\qquad
 \omega_b(x,h)=\sum_{r\in H}b(r^{-1}x)b(r^{-1}h^{-1}x).
 \tag{SM.18.3}
\]
Let measurable subsets $G_N$ increase to $G$, set
$\alpha_N=1_{G_N}$, and define on $D$
\[
 \beta_N(x)=\sum_{h\in H}\omega_b(x,h)
                       (1-\alpha_N(h^{-1}x)).
 \tag{SM.18.4}
\]
Then $0\leq\beta_N\leq1$, $\beta_N\to0$, and the omitted
absolute diagonal mass is bounded, uniformly for $|\rho|\leq1$, by
\[
 \int_D\sum_h|\rho(\pi x)d(\pi x)\omega_b(x,h)|
                 (1-\alpha_N(h^{-1}x))\,d\mu_D(x)
 \leq \int_D|d(\pi x)|\beta_N(x)\,d\mu_D(x)\longrightarrow0.
 \tag{SM.18.5}
\]
The initial-stratum and successive-shell diagonal densities
are absolutely summable in $L^1$, with sum of their norms equal
to $\|d\|_1$.  Thus diagonal integration and bounded cutoff
limits commute with this shell sum.  This conclusion does not
identify the pieces with independently specified local operators.

\begin{proof}
Write a singular-value expansion
$S=\sum_j\sigma_j u_j\otimes v_j^*$ with unit vectors $u_j,v_j$.
The series $\sum_j\sigma_j u_j\overline{v_j}$ converges in $L^1$,
since $\|u_j\overline{v_j}\|_1\leq1$ and
$\sum_j\sigma_j=\|S\|_1$.  Its sum is $d_S^\nu$.
The rank-one trace identity, followed by trace-norm and $L^1$
limits, proves (SM.18.1).  Testing against all bounded functions
gives uniqueness; the same estimate proves contractivity of
$S\mapsto d_S^\nu$.  Change of measure gives $d_S^\mu$.
The defining difference for the geometric kernel, with $\phi(1)=0$,
gives (SM.18.2), and dominated convergence gives the cutoff limit.

The discrete subgroup $H$ is countable.  All terms in $\omega_b$
are nonnegative, and Tonelli gives
\[
 \sum_h\omega_b(x,h)
 =\sum_r b(r^{-1}x)\sum_h b(r^{-1}h^{-1}x)=1.
 \tag{SM.18.6}
\]
Consequently $\beta_N$ lies in $[0,1]$.  For fixed $x,h$,
$\alpha_N(h^{-1}x)\to1$; dominated convergence with the
probability weights $\omega_b(x,h)$ gives $\beta_N(x)\to0$.
Another dominated convergence argument against $|d|$ proves
(SM.18.5), including uniformity in $\rho$.

Explicitly, let $c_0=\sum_h\omega_b(x,h)\alpha_0(h^{-1}x)$
and $c_j=\sum_h\omega_b(x,h)(\alpha_j-\alpha_{j-1})(h^{-1}x)$
for $j\geq1$.  Then $c_j\geq0$ and $\sum_jc_j=1$.
The shell densities are $d(\pi x)c_j(x)$, and Tonelli yields
$\sum_j\|dc_j\|_1=\|d\|_1$.  This proves every stated interchange.
\end{proof}

For a common-Hilbert-space realization of a proposed trace-class
operator remainder $E_N$, an equivalent trace-norm target is
\[
 \|d_{E_N}\|_1\leq\|E_N\|_1
   =\inf_{E_N=U_NV_N}\|U_N\|_2\|V_N\|_2\longrightarrow0.
 \tag{SM.18.7}
\]
The infimum runs over Hilbert--Schmidt factorizations.
The inequality $\|UV\|_1\leq\|U\|_2\|V\|_2$ is
Cauchy--Schwarz with trace duality.  Conversely, for
$E=W|E|$, take $U=W|E|^{1/2}$ and $V=|E|^{1/2}$;
both squared Hilbert--Schmidt norms equal $\operatorname{Tr}|E|$.
This proves the infimum formula.  Its vanishing is equivalent
to trace-norm convergence, and is sufficient for diagonal
convergence; it is not necessary for a scalar trace equality alone.
Hilbert--Schmidt convergence alone is insufficient: $N^{-1}$
times a rank-$N$ orthogonal projection has Hilbert--Schmidt norm
$N^{-1/2}$, trace norm one, and trace one.

Theorem~11.1.2.1, pp.~480--483, states a cofinal intersection
of closed carving subspaces and an identification by special
theta functions.  It does not state a common-space realization
and remainder estimate (SM.18.7).  On the other hand,
(SM.18.5) solves diagonal summation once the actual spectral
trace-class input and a normalized indentation are available.
Identifying those diagonal pieces with the local kernels remains
an additional assertion in the unwritten Section~12.3.

\subsection{The displayed orbital indentation vanishes}

\paragraph{Finding F-SM-19 (literal formula defect and nonzero test).}
Interpret equation~(11.3.2.2), p.~495, as the scalar-valued
function it is called, with the ordinary Haar measure
$d\mu_1(g)$ of its adjacent equations:
\[
 U(x,y)=\sum_{\xi\in F^\times}\int_{E_F^1}
   1_{\tau(\mathbb R_{>0})}(\xi gx)
   1_{\tau(\mathbb R_{>0})}(\xi gy)\,d\mu_1(g).
 \tag{SM.19.1}
\]
For $F=\mathbb Q$ this is identically zero, is not a normalized
indentation, and cannot yield the stated orbital trace formula
(11.3.2.1) under this literal interpretation.

\begin{proof}
Definition~A.2.3.6, p.~549, equation~(A.2.3.14), defines
$E_F^1$ as an actual subset of $\mathbb A_F^{\times,1}$.
For $\mathbb Q$ it is
$E_{\mathbb Q}^1=\{1\}_\infty\times\widehat{\mathbb Z}^{\,\times}$.
Definitions~A.3.2.1--2, pp.~559--560, use ordinary Haar measure
here, of volume one, with counting measure assigned only to
$\mathbb Q^\times$.  This infinite compact group has atomless
Haar measure: a positive singleton mass $c$ would give mass
$Nc$ to $N$ distinct points for every $N$, a contradiction.

Fix $x$ and $\xi$.  If $\xi gx=\tau(t)$ for $g\in E_{\mathbb Q}^1$,
the product formula and $|\tau(t)|=t$ force $t=|x|$.
Thus $g=\xi^{-1}\tau(|x|)x^{-1}$ is the only possible value.
The first indicator in (SM.19.1) is supported on at most one
point of $E_{\mathbb Q}^1$.  Each integral is zero and the
countable sum is zero, for all $x,y$.

For a nonzero right side of the orbital formula, choose
$F=\mathbb Q$, $\chi=1$, $S=\{\infty\}$, $a=1/2$, $s=4$,
and any $\alpha>0$.  Then $\operatorname{Pic}(\mathbb Z)=1$
and all stated ramification assumptions hold.
Choose $h\in C_c^\infty((3/2,5/2))$ with $h(2)=1$ and set
$\phi(u,t)=h(t)$ on
$C_{\mathbb Q}\simeq\widehat{\mathbb Z}^{\,\times}\times\mathbb R_{>0}$.
This is a Bruhat test function, finite-unit invariant and
supported over norms greater than one.  For the finite uniformizer
idele $\pi_2$, $|\pi_2|=1/2$, whence
\[
 \sum_{m\in\mathbb Z}\phi(\pi_2^m)
       =\sum_mh(2^{-m})=1.
\]
For the shell containing the place $2$, the stated right side is
$\log2/\mu_2(\mathbb Z_2^\times)>0$; the stated normalization
makes this $\log2$.  The literal indented kernel on the left
is zero, including after cutoffs, smoothening, and any integration
in the remaining norm variable.  The claimed equality therefore
fails with formula (SM.19.1) interpreted as printed.
\end{proof}

The introductory discussion, pp.~21--22, explicitly speaks of
scalar-valued measurable functions and allows indicators of
fundamental domains.  It introduces no Dirac interpretation
for the indicators in (SM.19.1).  A distributional replacement
or independently normalized disintegration would be a new
formula requiring its own operator and diagonal estimates.
This finding does not rule out such a corrected realization.

\subsection{A valid rational indentation and compact averaging}

\paragraph{Finding F-SM-20 (two explicit repairs and the remaining identification).}
For $\mathbb Q$, put
$Z=\mathbb R_{>0}\times\widehat{\mathbb Z}^{\,\times}
\subset\mathbb A_{\mathbb Q}^\times$ and $b=1_Z$.
Then $b$ is a continuous normalized indentation and
$b^\Delta(x,y)=\sum_\xi b(\xi^{-1}x)b(\xi^{-1}y)$ is valid.
For $x\in Z$ and $q\in\mathbb Q^\times$,
\[
 b^\Delta(x,qx)=1_{\{q=1\}}.
 \tag{SM.20.1}
\]
Hence this lift preserves the global diagonal but does not itself
recover the nontrivial prime-orbital formulas.

There is also a legitimate compact averaging for the actual operators.
Write $K=\widehat{\mathbb Z}^{\,\times}$ with probability Haar
measure, and assume $\phi(u,t)=h(t)$, $\rho(u,t)=\rho(t)$.
For $H_0=L^2(\mathbb R_{>0},w(t)\,dt/t)$, put
\[
 (Jf)(u,t)=f(t),\quad (J^*F)(t)=\int_KF(u,t)\,du,\quad E=JJ^*.
\]
Then $E$ is the orthogonal projection onto $K$-invariants, and
\[
 T_{\rho,\phi}=ET_{\rho,\phi}E,\qquad S_\phi=ES_\phi E.
 \tag{SM.20.2}
\]
For either operator $T$ in (SM.20.2), if $T$ is trace class,
\[
 \operatorname{Tr}_H T=\operatorname{Tr}_{H_0}(J^*TJ),
 \qquad \|T\|_1=\|J^*TJ\|_1.
 \tag{SM.20.3}
\]
If $K_T$ is its weighted Hilbert--Schmidt kernel, the compressed
kernel is the actual compact average
\[
 K_0(t,s)=\int_K\int_KK_T((u,t),(v,s))\,du\,dv,\qquad
 \|K_0\|_2\leq\|K_T\|_2.
 \tag{SM.20.4}
\]

\begin{proof}
For an idele $x$, let $r=\prod_p p^{v_p(x_p)}>0$, a finite
product, and take $q=\operatorname{sgn}(x_\infty)r$.
Then $q^{-1}x\in Z$.  Two such rational multipliers differ
by a rational number with all finite valuations zero and
positive sign, hence by one.  Thus the translates of $Z$
are disjoint and cover the ideles.  This subset is open
and closed and maps homeomorphically onto $C_{\mathbb Q}$.
Its intersection with the inverse image of a compact subset
of the quotient is compact.  Therefore $b$ has the stated
indentation properties and $\sum_qb(q^{-1}x)=1$.
For $x\in Z$ the first factor in $b^\Delta(x,qx)$ forces
$\xi=1$, and the second then forces $q=1$, proving (SM.20.1).

In particular, for the initial subgroup
$G_0=\mathbb R^\times\times\widehat{\mathbb Z}^{\,\times}$,
all outside-shell diagonal pieces from this lift are zero:
the only surviving $q$ is one and $x\in Z\subset G_0$.
The part labelled $S$ still contains the original global
spectral diagonal.  Its identification with a separately
defined $S$-local Paley--Wiener operator has not followed.

For averaging, Fubini and the volume-one normalization give
the displayed adjoint and $J^*J=1$, so $E$ is the claimed
orthogonal projection.  The actual weight $w_{I,s}$ depends
only on the norm.  Thus $K$ acts unitarily.
The closed geometric theta-function subspace of
Definition~10.4.1.1 is $C_F$-invariant, as explicitly used
in Lemma~10.4.1.2, p.~456.  That subspace and its orthogonal
complement are therefore $K$-invariant; their projections
commute with every $K$-translation and hence with $E$.
Convolution by $\phi$ satisfies
$EA(\phi)=A(\phi)=A(\phi)E$, by integrating the compact
variable, and $M_\rho$ commutes with $E$.  This proves
(SM.20.2).

Relative to $H=JH_0\oplus(JH_0)^\perp$, such a $T$ is
$(J^*TJ)\oplus0$.  The traces and singular values of this
block operator prove (SM.20.3).  Fubini first for simple
tensors and then $L^2$ density proves (SM.20.4); Jensen's
inequality for probability Haar measure on $K\times K$
gives its norm bound.
\end{proof}

These repairs permit a genuine trace calculation on the norm
coordinate for the indicated invariant part.  They do not
identify the compressed operator with the $S$-local operator
in Theorem~12.4.1.1, pp.~514--519.  That needs a specified
intertwiner carrying the actual kernels, projections and cutoff
limits to those local operators.  The zero indentation in
(SM.19.1) supplies none, and the closed-subspace carving
statement supplies no trace-norm estimate for that identification.

\subsection{Two precise obstructions to naive replacements}

\paragraph{Finding F-SM-21 (the unchanged archimedean local kernel misses a nonzero prime term).}
There exists a real $h\in C_c^\infty((3/2,5/2))$ such that
$h(2)=1$ and
\[
 P(h):=\widehat h(1/2)+\widehat h(-1/2)=0,\qquad A_0(h)=0,
 \tag{SM.21.1}
\]
where $A_0$ is the real, even archimedean distribution in F-SM-13.
Under the adopted Dirichlet theorem, its full zero distribution is
\[
 -\sum_\rho m(\rho)\widehat h(\rho-1/2)
       =\frac{\log2}{\sqrt2}>0.
 \tag{SM.21.2}
\]
Consequently the following two manuscript identifications cannot
be joined by simply replacing the zero indentation with $1_Z$
and leaving the $S=\{\infty\}$ local operator unchanged:
the global spectral trace identification of
Theorem~10.4.2.8 and equation~(12.5.1.4), and the local
Paley--Wiener trace limit of Theorem~12.4.1.1.
Under these explicitly stated operator inputs, the former
geometric cutoff trace tends to $\log2/\sqrt2$, whereas the
latter trace tends to zero, including the combined removed
moment correction $P(h)$.

More precisely, if trace-preserving transports put those
trace-class cutoffs on a common Hilbert space, their difference
$E_\lambda$ must satisfy
\[
 \liminf_{\lambda\to\infty}\|E_\lambda\|_1
       \geq \frac{\log2}{\sqrt2}.
 \tag{SM.21.3}
\]
This rules out that unchanged-kernel identification; it does not
rule out an assembly with the required nonzero prime contribution.

\begin{proof}
On the indicated interval, write
\[
 P(h)=\int h(t)p(t)\,dt,\quad
 A_0(h)=\int h(t)a(t)\,dt,\quad
 p(t)=\frac{t+1}{t^{3/2}},\quad
 a(t)=\frac{\sqrt t}{t^2-1}.
\]
The ratio
$R(t)=a(t)/p(t)=t^2/((t-1)(t+1)^2)$ is strictly decreasing for
$t>1$: its logarithmic derivative is
$(-t^2+t-2)/(t(t^2-1))<0$.
Choose nonzero nonnegative bumps $v_1,v_2$ supported,
respectively, in $(1.55,1.7)$ and $(2.3,2.45)$.
The quotients $A_0(v_i)/P(v_i)$ are weighted averages of
$R$ over these disjoint ordered intervals, so the first
is strictly larger than the second.  The matrix with
columns $(P(v_i),A_0(v_i))$ is therefore invertible.
Choose a real bump $u$ supported in $(1.9,2.1)$ with
$u(2)=1$, solve that two-by-two system for real $c_1,c_2$,
and set $h=u-c_1v_1-c_2v_2$.  This proves (SM.21.1) and
$h(2)=1$.

The only nonzero term in the finite distribution of F-SM-13 is
$\Lambda(2)h(2)/\sqrt2=\log2/\sqrt2$; the reciprocal terms
are outside the support.  Its explicit formula reads
$A_0(h)+B_1(h)=P(h)-\sum_\rho m(\rho)\widehat h(\rho-1/2)$,
so (SM.21.2) follows.  All zero sums converge absolutely.
The adopted zero-location input makes the same statement
hold in every positive-width strip used by the manuscript.

The source's global operator identification is exactly the
left side of (SM.21.2).  For $S=\{\infty\}$, its stated
local trace limit is $A_0(h)$, since $h(1)=0$ removes the
cutoff logarithm.  The combined moment correction also
vanishes by (SM.21.1).  Thus their trace difference tends
to $\log2/\sqrt2$.  The elementary inequality
$|\operatorname{Tr}E_\lambda|\leq\|E_\lambda\|_1$
proves (SM.21.3).
\end{proof}

\paragraph{Finding F-SM-22 (multiplying by a compact diagonal Dirac kernel destroys compactness).}
For the actual compact-invariant operators of F-SM-20, write
$T_{\rho,\phi}=E_K\otimes B_\rho$ on
$L^2(K)\otimes H_0$, where $E_K$ projects onto constants
and $B_\rho=J^*T_{\rho,\phi}J$.
If its radial kernel is $K_0(t,s)$, the proposed
multiplicative replacement
\[
 K_0(t,s)\longmapsto \delta_K(uv^{-1})K_0(t,s)
 \tag{SM.22.1}
\]
defines $1_{L^2(K)}\otimes B_\rho$ on tensor test functions.
If $B_\rho\ne0$, this operator is not compact and hence
is not trace class.  Under the operator inputs of F-SM-21,
$B_\rho\ne0$ for all sufficiently advanced cutoffs in that
example.  Thus (SM.22.1) cannot serve as a trace-class repair
of the ordinary-function formula.

\begin{proof}
The normalized compact diagonal distribution acts by
$\int_K\delta_K(uv^{-1})f(v)\,dv=f(u)$, so integrating
the product kernel gives $1\otimes B_\rho$ on simple
tensors and then on the Hilbert completion.
The group $K=\widehat{\mathbb Z}^{\,\times}$ has infinitely
many distinct continuous characters, forming an infinite
orthonormal set $(\chi_j)$ in $L^2(K)$.
Choose a unit vector $v\in H_0$ with $B_\rho v\ne0$.
The unit vectors $\chi_j\otimes v$ have images
$\chi_j\otimes B_\rho v$, which are mutually orthogonal
and have the same positive norm.  Their images have no
convergent subsequence.  The operator is not compact.
For the example in F-SM-21, its asserted global trace
limit is positive, so the trace-class $B_\rho$ has
nonzero trace and is nonzero for all sufficiently
advanced cutoffs.  This proves the final assertion.
\end{proof}

This is an obstruction to the particular \emph{multiplicative}
Dirac interpretation.  A subsequent compact-isotypic
projection changes the operator again: compression by $E_K$
restores $E_K\otimes B_\rho$.  Such a repair is possible
as an operation, but it returns to the global operator
and still does not supply the missing local identification.

\paragraph{Primary source and verification scope.}
Masullo, A. (2026). \emph{The music of the primes}
[supplied manuscript].
\url{https://math216.wordpress.com/wp-content/uploads/2026/11/the-music-of-the-primes.pdf}.
Additional exact locators: pp.~21--22; Lemma~12.2.1.1,
pp.~504--505; Proposition~12.2.2.2, pp.~508--510;
the missing Section~12.3, p.~514;
Theorem~11.3.2.1 and equation~(11.3.2.2), p.~495;
Definition~A.2.3.6, p.~549; Definitions~A.3.2.1--2,
pp.~559--560.  The p.~495 formula was inspected visually.
F-SM-18--22 prove the statements written here, without
certifying the earlier full trace-class construction or the
absent local-operator identification.
